% !PRINTED 98
\weberpage{98}
%<<CONT

where in the differentiation no account is taken of the dependence $a_i^{(k)}=a_k^{(i)}$. Then formula (9) becomes
\begin{equation}
 A\frac{\partial^2A}{\partial a_1^{(1)}\partial a_2^{(2)}}
 =\frac{\partial A}{\partial a_1^{(1)}}\frac{\partial A}{\partial a_2^{(2)}}
 -\left(\frac{\partial A}{\partial a_1^{(2)}}\right)^2.\tag{10}
\end{equation}

\begin{center}
{\large \S.\ 29.}

\vspace{0.5em}
{\large\gesperrt{Interpolation.}}
\end{center}

In § 9, under a special assumption, we already solved the problem of determining an integral rational function which assumes prescribed values for a sufficient number of prescribed values of the argument. We shall now solve this problem generally by applying determinants. First, however, we shall put forward the following observations.

In § 22, (12), the value of the determinant
\[
\begin{vmatrix}
 \alpha_1^{n-1}&\alpha_1^{n-2}&\cdots&\alpha_1&1\\
 \alpha_2^{n-1}&\alpha_2^{n-2}&\cdots&\alpha_2&1\\
 \cdots&\cdots&&\cdots&\cdots\\
 \alpha_n^{n-1}&\alpha_n^{n-2}&\cdots&\alpha_n&1
\end{vmatrix}
\]
was found to be equal to the product of differences
\begin{equation}
(\alpha_1-\alpha_2)(\alpha_1-\alpha_3)\cdots(\alpha_1-\alpha_n)
(\alpha_2-\alpha_3)\cdots(\alpha_2-\alpha_n)\cdots(\alpha_{n-1}-\alpha_n).\tag{1}
\end{equation}
For brevity we now denote this product of differences by
\begin{equation}
 [\alpha_1,\alpha_2,\alpha_3,\ldots,\alpha_n],\tag{2}
\end{equation}
so that the quantity $[\alpha_1,\alpha_2,\alpha_3,\ldots,\alpha_n]$ changes its sign when two of the elements $\alpha_1,\alpha_2,
\ldots,\alpha_n$ are interchanged.

We further define an integral rational function $f(x)$ of degree $n$ in the variable $x$ by the equation
\begin{equation}
 f(x)=(x-\alpha_1)(x-\alpha_2)\cdots(x-\alpha_n),\tag{3}
\end{equation}
so that, by § 13,
\begin{align}
 f'(\alpha_1)&=(\alpha_1-\alpha_2)(\alpha_1-\alpha_3)\cdots(\alpha_1-\alpha_n),\notag\\
 f'(\alpha_2)&=(\alpha_2-\alpha_1)(\alpha_2-\alpha_3)\cdots(\alpha_2-\alpha_n),\tag{4}\\
 &\vdots\notag\\
 f'(\alpha_n)&=(\alpha_n-\alpha_1)(\alpha_n-\alpha_2)\cdots(\alpha_n-\alpha_{n-1}).\notag
\end{align}
%CONT>>
